% Research module prepared for incorporation into DELTA_ON_GT_paper.tex.
% Local statement labels LP and LB are deliberately reserved for root renumbering.
% Attribution: AI-produced by GPT-6 Astra (OpenAI, ChatGPT Pro mode).

\subsection{Integral realization and the strict profinite inclusion}

\paragraph{LP.1 (Proved): integral peripheral realization.}
Let $P_\ell$ be free pro-$\ell$ on $x,y$, let
$z=(xy)^{-1}$, and let $L=\operatorname{Lie}_{\mathbf Z_\ell}(X,Y)$.
Write $\gamma_m$ for the closed ordinary lower central series. Its
associated graded Lie algebra is $L$, by the integral Magnus map.
For $n\geq2$, every homogeneous special outer derivation of degree $n$
is the leading class of a continuous multiplier-one peripheral outer
automorphism of $P_\ell$. This statement holds at every prime.

\emph{Proof.}
We first justify the two properties of the lower central comparison which
we need. The Magnus expansion sends the generators to $1+X,1+Y$;
commutators in $\gamma_m$ have leading degree $m$, and the leading
terms of the standard Lyndon commutators give the integral free-Lie
basis. Consequently the comparison is an isomorphism on every
$\gamma_m/\gamma_{m+1}$, with coefficients in $\mathbf Z_\ell$.
It is separated and complete for this filtration: every finite
$\ell$-group is nilpotent, so the intersection of the $\gamma_m$ is
trivial, and compactness identifies $P_\ell$ with
$\varprojlim_m P_\ell/\gamma_m$.

Choose a representative of the special outer derivation fixing $X$.
Its two other peripheral conjugators are homogeneous $h,k\in L_n$,
with
\[
             [h,Y]=[k,X+Y].
\]
Choose actual elements of $\gamma_nP_\ell$ with these leading terms.
Their provisional images of $y,z$, together with $x$, satisfy the
peripheral product identity through degree $n+1$. If the error is in
$\gamma_m$, alter the two conjugators in degree $m-1$. The change
in the degree-$m$ error is
\[
 L_{m-1}\oplus L_{m-1}\longrightarrow L_m,
 \qquad (u,v)\longmapsto[u,Y]+[v,-X-Y].
\]
This map is onto integrally. Its image is
$[L_{m-1},X]+[L_{m-1},Y]$, and Jacobi repeatedly moves a generator
to the outside of any bracket of total degree $m$. Thus every error
can be killed without modifying the prescribed leading terms.
The corrections converge in every finite $\ell$-group quotient,
and therefore in $P_\ell$. They give an exact peripheral endomorphism.
It induces the identity on the Frattini quotient, so the pro-$\ell$
Burnside basis theorem makes it surjective.

For completeness, the Hopficity step is as follows. If $\alpha:G\to G$
is a continuous epimorphism of a finitely generated profinite group,
then $N\mapsto\alpha^{-1}N$ injects the finite set of open normal
subgroups of a given index into itself, hence permutes it. For any
open normal $N$, some $k>0$ has $\alpha^{-k}N=N$. Since
$\ker\alpha\subseteq\ker\alpha^k\subseteq\alpha^{-k}N=N$,
the kernel lies in every open normal subgroup and is trivial.
Applying this to the endomorphism above proves the assertion.
Inner derivations are realized by inner automorphisms, so the
construction passes to outer classes. No exponential or division by
factorials is used. \hfill$\square$

\paragraph{LP.2 (Proved): coprime geometric symmetry lifts.}
For $\ell\geq5$, every integral special outer leading class fixed by
the geometric $S_3$ lifts to an actual special outer automorphism
centralizing that $S_3$.

\emph{Proof.}
Let $I^n$ be the closed subgroup of IA special outer automorphisms
whose first possible derivation degree is $n$. It is pro-$\ell$:
in each finite characteristic $\ell$-group quotient, the automorphisms
trivial on the Frattini quotient form an $\ell$-group. The geometric
$S_3$ acts by conjugation on $I^n$ and preserves $I^{n+1}$.
LP.1 proves surjectivity onto the full integral special leading module.
Given a lift $a$ of an invariant class, the elements $a^{-1}g(a)$,
$g\in S_3$, form a nonabelian $1$-cocycle in $I^{n+1}$.
Its class is zero: in a finite quotient this is the conjugacy of
complements in an $\ell$-group-by-$S_3$ extension, since $6$ is
prime to $\ell$. One may prove this by induction on an invariant
central elementary abelian subgroup, averaging the additive
cocycles there by $1/6$. Compactness of the finite solution sets
gives a single correction in $I^{n+1}$ at the inverse limit.
Correcting $a$ makes it $S_3$-fixed without changing its leading
term. This is a group construction, not a claim that an arbitrary
graded solution exponentiates integrally. \hfill$\square$

\paragraph{LP.3 (Proved): exact normalized pro-$\ell$ pairs.}
The centralizer in LP.2 is identified with the multiplier-one
two-cycle and three-cycle pairs on $P_\ell$, with their actual
automorphism condition. If its first outer derivation is $-[D_h]$,
where $D_h(X)=0$ and $D_h(Y)=[h,Y]$ has degree $n\geq2$, the
normalized pair $f_\ell$ lies in $\gamma_nP_\ell$ and has leading
term exactly $h$.

\emph{Proof.}
Put $E_\ell=\widehat B/O^\ell(F)$. It is the actual relative
completion in
\[
 1\longrightarrow P_\ell\longrightarrow E_\ell
             \longrightarrow S_3\longrightarrow1,
 \qquad x=(sr)^{-2},\quad y=sxs.
\]
The center of $P_\ell$ is trivial, and $S_3$ acts faithfully on
$P_\ell^{\rm ab}$. Thus conjugation embeds $E_\ell$ into
$\operatorname{Aut}(P_\ell)$ and identifies it with the pullback
of $\operatorname{Aut}(P_\ell)\to\operatorname{Out}(P_\ell)$ along
the geometric map $S_3\to\operatorname{Out}(P_\ell)$.
An outer automorphism centralizing this map therefore extends,
uniquely up to an inner automorphism by $P_\ell$, to $E_\ell$,
inducing the identity on $S_3$.

Choose a special representative fixing $x$. Its extension sends
$\tau=sr$ to an element with the same square $x^{-1}$ and the
same image in $S_3$. Their quotient lies in $C_{P_\ell}(x)=\langle
x\rangle$. If it is $x^v$, equality of squares gives $2v=0$ in
$\mathbf Z_\ell$, hence $v=0$. Thus $\tau$ is fixed literally.
The centralizer assertion follows from the Magnus leading terms:
after removing the abelian power of $x$, a first nonzero term
commuting with $X$ would belong to the centralizer of $X$ in
$L_m$, which is zero for $m>1$. That last assertion follows by
embedding in the free associative algebra; a homogeneous tensor
commuting with $X$ is a multiple of $X^m$, and no positive
commutator Lie element has such a term.

By coprime complement conjugacy, the new $s,r$ are conjugate to
the old positive generators by elements of $P_\ell$. Put $f=s\Phi(s)$.
Since $\Phi(s)\Phi(r)=sr$, we have
\[
 \Phi(s)=sf,\qquad \Phi(r)=f^{-1}r,
 \qquad \Phi(y)=f^{-1}yf.
\]
The first torsion relation makes $f^{\rm ab}=(a,-a)$. Replacing
$\Phi$ by $\operatorname{Inn}(x^a)\Phi$ keeps $\tau$ fixed and
replaces $f$ by $y^af x^{-a}\in P_\ell'$. In the marked group
conjugation by $s$ exchanges $x,y$, and conjugation by $r$
cycles $x,y,z$, where $z=(xy)^{-1}$. The two torsion relations
therefore give, with this exact order,
\[
 f(y,x)f(x,y)=1,\qquad f(z,x)f(y,z)f(x,y)=1.
\]
The restriction of $\Phi$ is already an automorphism. Conversely,
these equations and that automorphism condition define an
automorphism of $E_\ell$ by the displayed assignments. Its outer
restriction centralizes $S_3$ because the quotient action is the
identity. This also gives a direct proof of the pro-$\ell$ version
of Harbater--Schneps, Propositions 3--4 and Theorem 2, rather than
silently transporting their profinite result to a different category.

It remains to verify the leading coefficient. In
$P_\ell/\gamma_{n+1}$, the centralizer of $x$ is
$\langle x\rangle\gamma_n/\gamma_{n+1}$, by the same leading-term
argument. If $\psi_f$ is inner there, say by $w$, then $w$ belongs
to that centralizer and $fw$ belongs to the corresponding centralizer
of $y$. Since $f^{\rm ab}=0$, both powers vanish and
$f\in\gamma_n$. If its leading term is $h_f$, then
$\psi_f(y)y^{-1}$ has leading term $-[h_f,Y]$.
Equality of the prescribed outer class with $-[D_h]$, followed by
the free-Lie centralizer argument at $X$ and then at $Y$, gives
$h_f=h$. All maps here are continuous quotient or conjugation maps.
\hfill$\square$

\paragraph{LP.4 (Proved): an independently checked nonpentagonal leading term.}
For the standard Lyndon brackets $b_w$ with $x<y$, set
\[
\begin{split}
q_7={}&3b_{xxxyxyy}+3b_{xxxyyxy}+b_{xxyxyxy}
       -3b_{xxyxyyy}\\
     &-3b_{xxyyxyy}-3b_{xxyyyxy}-b_{xyxyxyy}.
\end{split}
\]
It has integral coefficients, is nonzero, and satisfies
\[
 q_7(x,y)+q_7(y,x)=0,\qquad
 [x,q_7(x,y)]=[x+y,q_7(-x-y,y)].
\]
It is special and has zero three-cycle norm. Its infinitesimal
pentagon defect is nonzero, already over every field of
characteristic at least five.

\emph{Proof and exact certificate.}
The program \texttt{check\_lie\_witnesses.py} expands standard
Lyndon brackets by splitting at the longest proper Lyndon suffix.
All four asserted identities are checked coefficientwise over the
integers; $q_7$ has 62 nonzero associative coefficients. Here are the
inner terms needed for the symmetry assertion. Put
$\alpha(X)=Y$, $\alpha(Y)=X$, and
$\theta(X)=-X-Y$, $\theta(Y)=Y$. They are the transpositions
$(X,Y)$ and $(X,Z)$ on $X+Y+Z=0$ and generate $S_3$. If
$D_h(X)=0$, $D_h(Y)=[h,Y]$ and $k=\theta(h)$, antisymmetry and
the displayed cabling equation give exactly
\[
 \alpha D_h\alpha^{-1}-D_h=\operatorname{ad}(-h),\qquad
 \theta D_h\theta^{-1}-D_h=\operatorname{ad}(k-h).
\]
For example the latter difference sends $X$ to $-[k,Y]$ and
$Y$ to $[k-h,Y]$, while cabling says
$[k-h,X]=-[k,Y]$. Thus $[D_h]$ is fixed by both generators.
The program additionally checks these two derivation identities on
both free generators for $q_7$ and $h_9$.

We retain an actual braid comparison. The kernel of
$P_4\to P_3$ is free on the three braids $A_{14},A_{24},A_{34}$.
Conjugation by $P_4$ on this kernel, followed by its pro-$\ell$
completion, is an actual continuous group action; the full twist
acts trivially, so it factors through the pure five-point sphere
group. Its leading derivations on $L(a,b,c)$ are
\[
 t_{12}\mapsto D_{ab},\quad t_{13}\mapsto D_{ac},\quad
 t_{23}\mapsto D_{bc},\qquad
 t_{14}\mapsto\operatorname{ad}a,\quad
 t_{24}\mapsto\operatorname{ad}b,\quad
 t_{34}\mapsto\operatorname{ad}c,
\]
where $D_{uv}(u)=[u,v]$, $D_{uv}(v)=[v,u]$, and $D_{uv}$
kills the third generator. These formulas follow by linearizing
the conjugation presentation in Cohen, equations (1.2)--(1.4).
They are therefore leading terms of the actual group action,
not merely a representation of an unrelated presented Lie algebra.
The program independently checks all 18 required infinitesimal
braid relations on these derivations.

For the pentagon convention
\[
\begin{split}
\mathcal P(h)={}&h(t_{12},t_{23}+t_{24})
 +h(t_{13}+t_{23},t_{34})-h(t_{23},t_{34})\\
 &-h(t_{12}+t_{13},t_{24}+t_{34})-h(t_{12},t_{23}),
\end{split}
\]
the program obtains 880, 882 and 808 nonzero coefficients on
$a,b,c$, respectively. The coefficient of $a^5bcb$ in
$\mathcal P(q_7)(a)$ is exactly $16$. The complete 2,570-coefficient
output is retained in \texttt{lie\_witnesses\_exact.json}.
This coefficient proves nonvanishing over every $\mathbf Q_\ell$
and modulo every $\ell\geq5$. \hfill$\square$

\paragraph{LP.5 (Proved, using R1): strictness for the full profinite groups.}
There is an element of $\widehat{\mathrm{GT}}_0$ outside
$\widehat{\mathrm{GT}}$. In particular
\[
             \widehat{\mathrm{GT}}
                    \subsetneq\widehat{\mathrm{GT}}_0.
\]
This proves the corollary requested in P.2. It assigns no value to
$\Delta$ on $\widehat{\mathrm{GT}}$.

\emph{Proof.}
Fix any $\ell\geq5$. LP.1--LP.3 lift the invariant class
$-[D_{q_7}]$ to an actual normalized pair
$f_\ell\in\gamma_7P_\ell\cap P_\ell'$ whose leading term is $q_7$.
Its pure map is an automorphism and it satisfies both group cycle
relations. The last assertion of the refereed R1 applies to
precisely this pair, retaining its entire pro-$\ell$ image. It
gives $(1,f)\in\widehat{\mathrm{GT}}_0$ with
$q_\ell(f)=f_\ell$.

If this full pair satisfied the pentagon, apply the continuous
quotient to pro-$\ell$ pure braid groups and then the action on
the free pro-$\ell$ kernel used in LP.4. Taking the first
nonzero Johnson degree of the resulting identity gives
$\mathcal P(q_7)=0$. To justify this direction explicitly,
substitution of a group element whose first term has degree seven
depends in that degree only on the degree-one terms of its
arguments. Products of the five substituted elements add in that
degree; all commutators between them and all higher corrections
have larger degree. The cabling and quotient homomorphisms used
in the group pentagon induce exactly the sums of $t_{ij}$ in
LP.4. Thus a zero group defect forces a zero displayed leading
derivation. Its coefficient 16 contradicts that conclusion.
Every full lift of this same pro-$\ell$ pair is excluded for this
reason; no inference from a surviving finite shadow is involved.
\hfill$\square$

\subsection{The two weight calculations and the scope of metabelian lifting}

\paragraph{LB.1 (Proved, using R1): the full special weight-five image.}
For every prime $\ell$, the rational weight-five metabelian leading
image of the multiplier-one full peripheral outer automorphism
group is
\[
 (X+Y)\langle X^2,XY,Y^2\rangle_{\mathbf Q_\ell}[x,y].
\]
It has dimension three. Thus LM7\_PUSH D.51 is correct, and
its rational dimension statement also extends to two and three.

\emph{Proof.}
In the special equation $[h,y]=[k,x+y]$, take the actual quotient
of the free Lie algebra by its second derived ideal. If the
coefficients are $p,q$, then $Yp=(X+Y)q$. Since $Y,X+Y$ are
relatively prime, $X+Y$ divides the homogeneous cubic $p$.
This proves the upper bound three.

For the lower bound the three pairs in LM7\_PUSH D.51.3 were
expanded anew in \texttt{check\_lie\_witnesses.py}. All three
special equations vanish coefficientwise. Right-$y$ Fox
extraction, followed by abelianization, gives $XY$ times their
metabelian coefficient. The output is respectively
$(X+Y)X^2,(X+Y)XY,(X+Y)Y^2$. Their only possible denominator is
two. For odd $\ell$, LP.1 integrates the displayed vectors in
the actual free pro-$\ell$ group. At $\ell=2$, integrate twice
each displayed pair; their rational images still span the same
three-dimensional $\mathbf Q_2$-space. No coprime $S_3$ step is
needed for this unrestricted peripheral statement.
Place each automorphism in the $\ell$ factor of
$F/R_{\rm nil}(F)=\prod_pF^{(p)}$, and place the identity in all
other factors, with the same tuple of peripheral conjugators.
The free-group assertion of R1 lifts this exact automorphism.
The characteristic pro-$\ell$ projection and the lower-central
and metabelian quotient maps retain the three leading terms.
Their linear independence proves the lower bound.

The representation-theoretic upper bound in D.21 also follows:
division by $X+Y$ identifies the target with quadratics in the
standard two-dimensional $S_3$ representation. In characteristic
zero this is the invariant line $X^2+XY+Y^2$ plus an irreducible
two-dimensional summand. An $S_3$-stable image containing the
invariant line therefore has dimension one or three; the actual
source construction decides it is three. \hfill$\square$

\paragraph{LB.2 (Proved, using R1): the full gentle weight-nine image.}
For every $\ell\geq5$, the rational weight-nine metabelian leading
image of $\widehat{\mathrm{GT}}_0$ is
\[
 \mathbf Q_\ell G_9[x,y]\oplus\mathbf Q_\ell P_9[x,y],
 \qquad G_9=\frac{(X+Y)^9-X^9-Y^9}{XY},\quad
 P_9=X^2Y^2(X+Y)^3.
\]
There is an actual normalized full pair whose first pro-$\ell$
coefficient has metabelian image $P_9[x,y]$.

\emph{Proof.}
The 42 displayed coefficients of $2h$ in LM7\_PUSH D.52.3
were entered afresh and expanded over $\mathbf Q$. The
antisymmetry, cabling, special and cyclic equations all vanish
coefficientwise; the associative support has 412 terms. Its Fox
image is exactly $XY P_9$, with coefficient vector $(1,3,3,1)$
in the four indicated monomials. Hence this is an integral
invariant special source at every $\ell\geq5$. LP.1--LP.3 and
R1 give the actual full pair with this leading term and its
correct positive pair sign. The derivation coordinate is its
negative, because $y\mapsto f^{-1}yf$.

For the upper bound let $p$ be the homogeneous degree-seven
metabelian coefficient of an invariant special leading term.
Its equations include
\[
 p(X,Y)=p(Y,X),\qquad
 Xp(X,Y)+(X+Y)p(-X-Y,Y)=0,
\]
and its cyclic sum is zero. Comparing coefficients gives
\[
\begin{split}
p={}&A(X^7+Y^7)+4A(X^6Y+XY^6)\\
   &+B(X^5Y^2+X^2Y^5)+(3B-14A)(X^4Y^3+X^3Y^4)\\
 ={}&\frac A9G_9+\left(B-\frac{28A}{3}\right)P_9.
\end{split}
\]
Thus the upper bound is two. We also independently formed all
56 Lyndon basis vectors in degree nine, formed the 2002 by 56
integer matrix of the three identities, and computed its exact
rational kernel. It has dimension four and metabelian rank two.
An explicit invariant free-Lie preimage of $G_9$ with denominator
four is retained in \texttt{invariant\_dimensions.json}. LP.1--LP.3
and R1 integrate that vector as well. This gives both directions
without using a circular dependence on D.1 or D.2, or an unproved
Galois nonvanishing statement. The exact matrix and all kernel
vectors are reproducible from \texttt{check\_invariant\_dimensions.py}.
\hfill$\square$

\paragraph{LB.3 (Known with application proved): every weight-nine lift with this metabelian coefficient fails.}
Enriquez's Theorem 0.1 states that a homogeneous degree-$m$
element of $\mathfrak{grt}_1(k)$ has metabelian image proportional
to $G_m[x,y]$ for odd $m\geq3$, and zero in other degrees.
Here $k$ is any characteristic-zero field. Consequently every
full gentle pair with no pro-$\ell$ terms below degree nine and
with degree-nine metabelian image $P_9[x,y]$ fails the pentagon,
regardless of its other degree-nine or higher coefficients.

\emph{Proof of the application.}
If a full pair in this stratum satisfied the group pentagon, its
first nonzero term would satisfy the infinitesimal pentagon, by
the continuous comparison and first-degree argument of LP.5.
Its two group cycles likewise give the antisymmetry and cyclic
Lie identities. It therefore belongs to
$\mathfrak{grt}_1(\mathbf Q_\ell)_9$ and Enriquez applies.
The coefficient of $X^7$ in $G_9$ is 9, whereas it is zero in
$P_9$, and $P_9$ is nonzero. Thus they are not proportional.
The contradiction excludes every lift in this entire stratum.
This also supplies the requested pentagon upper bound of
LM7\_PUSH D.20; it does not by itself prove the separate
Galois lower-bound assertion in that row. \hfill$\square$

\subsection{The complete finite frame image and the second mixed pair}

\paragraph{LB.4 (Proved from the verified frame comparison): the finite image.}
In the signed corrected coordinates of R4,
\[
 (T,W):\mathcal H\longrightarrow C_6\times C_4
\]
is onto. The image of $G_{\mathbf Q^{\rm ab}}$ is precisely
\[
 J=\{(T,W):\epsilon_3(T)=\epsilon_2(W)\},
\]
which is cyclic of order 12.

\emph{Proof.}
Write $C_6=\langle -I\rangle\times\langle R_0\rangle$ and
$C_4=\langle S\rangle$, with the matrices and signed lift
which are identity modulo four as in R4. R4 gives the
homomorphisms and the containment of the Galois image in $J$.
Writing the two cyclic exponents as $i\bmod6,j\bmod4$ identifies
$J$ with $i\equiv j\pmod2$; it is generated by $(1,1)$ and has
order 12.

The three-division polynomial of $E_3:y^2=x^3+1$ is
$3x(x^3+4)$. The nonzero roots have $y^2=-3$, so
\[
 \mathbf Q(E_3[3])=\mathbf Q(\zeta_3,\sqrt[3]4).
\]
Its Galois group over $\mathbf Q$ is $S_3$, since $X^3-4$ is
irreducible and has nonsquare discriminant. Its intersection
with $\mathbf Q^{\rm ab}$ is the quadratic fixed field of
$A_3$, so the action over $\mathbf Q^{\rm ab}$ has order three.
The normalization here is essential. In the transported Legendre
basis used in R4's referee report, section 1, step 9.4, one has the
signed formula
\[
 T_\sigma=\kappa(\sigma)\rho_{E_3,12}(\sigma),\qquad
 \kappa(\sigma)=\sigma(r_0)/r_0\in\{\pm1\},
 \quad \sigma\in G_{\mathbf Q^{\rm ab}}.
\]
It follows by writing the equianharmonic Legendre fiber as the
quadratic twist of $y^2=x^3+1$ by $1/\sqrt{-3}$; the square root
has the same character as $r_0$, since $\sqrt2$ and $\zeta_8$
belong to $\mathbf Q^{\rm ab}$. Reduction modulo three and
quotienting by $\langle-I\rangle$ removes exactly this quadratic
twist. The resulting projective torsion representation is therefore
the projection of $T$ to $C_6/\langle-I\rangle$, with no further
unspecified normalization. Its image has order three, so the
Galois frame image has order divisible by three.

Let $r_0=(3/4)^{1/4}$ and $M=\mathbf Q(i,r_0)$, the field
identified in the refereed proof of R4. The polynomial
$X^4-12$ of $2r_0$ is Eisenstein at three, and
$\mathbf Q(r_0)$ is real. Therefore $[M:\mathbf Q]=8$, with
the dihedral action $r_0\mapsto i r_0$, $i\mapsto i$ and
complex conjugation. Its maximal abelian subfield is
$\mathbf Q(i,r_0^2)=\mathbf Q(i,\sqrt3)$. Hence an element of
$G_{\mathbf Q^{\rm ab}}$ sends $r_0$ to $-r_0$. The verified
sign convention of R4 makes both parities one for that element.
Its $W$ component is an odd power of $S$ and has order four.
The Galois image in the cyclic group $J$ thus has order divisible
by both three and four, so it is all of $J$.

For the element of R2, the corrected coordinates are $(-I,I)$.
It is outside $J$. Consequently the image of $\mathcal H$
contains the order-12 subgroup $J$ and its other coset, and is
all of $C_6\times C_4$. The group enumeration is independently
checked in \texttt{check\_cm\_replacement.py}. \hfill$\square$

\paragraph{LB.5 (Proved, using R2): the arithmetic replacement lifts.}
Let $N=\ker q_{12}\subset\widehat B$. There is a continuous
positive peripheral automorphism of
$\widehat B/R_{\rm nil}(N)$ with the LM7 D.46 component
description and specified conjugators. It lifts by R2 to an
element of $\mathcal H$ with projective pair $(I,[i])$, parities
$(0,1)$ and $\Delta=1$. We can prove the needed arithmetic
input without assuming LM7 D.35's entire 2-adic image theorem.

\emph{Proof.}
Put $k=\mathbf Q(\zeta_3)$ and
$a=1-2\zeta_3^2=3+2\zeta_3=2+\sqrt{-3}$. Counting
$E_3(\mathbf F_7)$ gives 12 points, so Frobenius at either
prime of $k$ over seven has characteristic polynomial
$X^2+4X+7$. The automorphism $(x,y)\mapsto(\zeta_3x,y)$
makes $T_2E_3$ a free rank-one module over
$\mathbf Z_2[\zeta_3]$. Frobenius commutes with this module
structure and is multiplication by one of $-a,-\bar a$.
Choose the prime, or conjugate it, so that it is $-a$.

All two-torsion of $E_3$ is $k$-rational. Its full 2-division
extension over $k$ is therefore pro-two: its image is in the
pro-two kernel of reduction modulo two. Its intersection with
$k(\sqrt[3]4)=k(E_3[3])$ is $k$, since the latter extension is
cyclic of degree three. The resulting product surjection of
Galois groups permits an actual $g\in G_k$ with 2-adic image
$-a$ and trivial action on $E_3[3]$. This uses only one actual
Frobenius and the coprime disjointness of these extensions.

Choose $\tau\in I_3$ inducing the nontrivial automorphism of
$k$, which is ramified at three, and put $h=[g,\tau]$.
The full cyclotomic character of a commutator is one. The
normalized action of $I_3$ on the free pro-two sphere group is
trivial by good reduction, with the integral tangential base
point $\overrightarrow{01}$. This exact automorphism statement
is stated in Hain--Matsumoto, section 8.5 immediately before
Theorem 8.7, for $N=1,\ell=2$. Consequently the normalized
action of $h$ on the entire pure pro-two group is the identity.
On $T_2E_3$, conjugation by $\tau$ conjugates the CM coefficient
ring. Thus
\[
 \rho_{E_3,2}(h)=a/\bar a=(5+8\zeta_3)/7
                    =(-3-8\zeta_3^2)/7.
\]
It is $-I$ modulo four. Also $h$ acts trivially on $E_3[3]$,
since $g$ does. The signed frame convention of R4, which is
$I$ modulo four, therefore gives $T_h=-I$ modulo three.
Its first parity is one and, by R4, its second parity is one.

We check that the induced action on $Q$ is literally the identity.
For every multiplier-one pair, the pure $r$-frame $A$ belongs
to $F'$. Indeed $f=rAr^{-1}A^{-1}$, so
$(r-1)A^{\rm ab}=0$; the action of $r-1$ on
$F^{\rm ab}=\widehat{\mathbf Z}^2$ is injective, with determinant
three. Hence $q_3(A)$ is in the quaternion derived subgroup of
$\mathrm{SL}_2(\mathbf F_3)$. In its abelianization $C_3$, the
identity $T=q_3(A)^{-1}U^{8\rho}$ gives $[T]=2\rho[U]$.
Since $T_h=-I$ has trivial class there, $\rho_h\equiv0\pmod3$.
The shear $U^{8\rho_h}$ is consequently $I$ modulo twelve.
R3 now gives $q_{12}\Phi_h=q_{12}$. Its frame coordinates in
$Q$ are $a_Q=[5I]$ and some $b_Q$ with inverse reducing to
$\pm S$ at three and to $I$ at four.

For each prime $p$, put $E_p=\widehat B/O^p(N)$. The actual
fiber product theorem of R2 gives
\[
 \widehat B/R_{\rm nil}(N)=\prod_Q E_p.
\]
Because $\Phi_h$ preserves $N$ and is the identity on $Q$,
the tuple with $\Phi_h$ in the two-component and identity in
all odd components is a continuous automorphism of this
fiber product fixing $\tau$. We must specify its peripheral
conjugators; a product of automorphisms alone would not suffice.

In $E_2\to Q$, the element $a_Q^{-1}$ centralizes $r_Q$.
Coprime complement conjugacy for the order-three action lifts
it to $c_2\in C_{E_2}(r)$. Therefore $A_{h,2}c_2$ is a
conjugator of $r$ with $Q$-coordinate one. At every odd $p$,
the element $b_Q$ centralizes $s_Q$ and lifts by the same
argument to $b_p\in C_{E_p}(s)$. The two actual tuples
\[
 A^\dagger=(A_{h,2}c_2,(1)_{p\ne2}),\qquad
 \beta^\dagger=(\beta_{h,2},(b_p)_{p\ne2})
\]
have common $Q$-coordinates $1,b_Q$. They belong to the image
of $F$ and are the required positive peripheral conjugators;
the cusp conjugator is one. Thus every hypothesis of R2 is
verified, including the prime-two and prime-three cases and
the mixed finite base $Q$.

R2 lifts the automorphism with these conjugators to $\widehat B$
and keeps $\tau$ fixed. The resulting parameter $f=s\Phi(s)$
is trivial in every absolute pure pro-$p$ quotient: for $p=2$
this was proved for $h$, and for odd $p$ the chosen component
is the identity. More explicitly, an identity pure action gives
$f_p\in C_{F^{(p)}}(y)=\langle y\rangle$, while the actual
torsion relation $(sf_p)^2=1$ makes its abelian coordinate
$(a,-a)$. An element of $\langle y\rangle$ has zero first
coordinate, so $a=0$ and this element is one.
In particular $f\in F'$. The torsion word
calculation in LP.3 now proves both cycles, and the restriction
is an actual automorphism. Hence it is an element of $\mathcal H$.
Its induced action on $Q$ is still literally the identity.
By R3 its shear $U^{8\rho}$ therefore centralizes the image of
$F$; modulo three that image is $\mathrm{SL}_2(\mathbf F_3)$.
The intersection of its center with the upper unipotent group
is trivial, so $\rho\equiv0\pmod3$ and the shear is $I$
modulo twelve. This distinguishes the corrected shear from the
separately specified cusp conjugator one.
Its retained $Q$-frames have $A_Q=1$, $\beta_Q=b_Q$, and
corrected $V_Q=1$, giving $T=I$, $W=\pm S$. Projectively this is the
pair $(I,[i])$ claimed in LM7 D.46, and its parities are $(0,1)$.

In the graph basis in which $\zeta_3^2$ acts by
$\left(\begin{smallmatrix}0&1\\-1&-1\end{smallmatrix}\right)$,
the computed scalar is
$\left(\begin{smallmatrix}11&8\\8&3\end{smallmatrix}\right)$
modulo sixteen. This and the Frobenius point count are recorded
by \texttt{check\_cm\_replacement.py}. The construction does
not assert that this element satisfies the pentagon.
\hfill$\square$

\paragraph{LB.6 (Proved): the discrepancy is independent of absolute pronilpotent data.}
Let $v:\widehat{\mathrm{GT}}_0\to V$ be the product of its
specified absolute pro-prime, pronilpotent and associated graded
actions, with $V$ its image. Then
\[
 (v,\Delta):\widehat{\mathrm{GT}}_0
                  \longrightarrow V\times\mathbf F_2
\]
is a continuous epimorphism. The Galois image has zero second
coordinate. With $V_{\mathcal H}=v(\mathcal H)$, the stronger
map $(v,\epsilon_3T,\epsilon_2W):\mathcal H\to
V_{\mathcal H}\times\mathbf F_2^2$ is onto.

\emph{Proof.}
If $U$ is any subgroup of $F$ and $q:F\to D$ has finite
nilpotent image, then $q(U)$ is nilpotent. Hence
$R_{\rm nil}(U)\subseteq R_{\rm nil}(F)$. The R2 witness,
which is identity on $\widehat B/R_{\rm nil}(F'F^3)$, is
therefore identity on every absolute pro-$p$ quotient of $F$,
on their product, and on all the associated graded quotients.
It lies in $\ker v$ and has $\Delta=1$. Multiplying any chosen
lift of a point of $V$ by this element toggles the second
coordinate. This proves the first epimorphism.

The construction in LB.5 is also identity on every absolute
pure pro-$p$ quotient and has parity pair $(0,1)$; the R2
witness has pair $(1,0)$. Both lie in $\ker v$. Their products
realize every parity pair in that kernel, proving the stronger
assertion on $\mathcal H$. The restriction on the Galois image
is exactly R4. \hfill$\square$

The word ``quotients'' in this assertion refers to the geometric
actions induced from the absolute pro-prime and pronilpotent
quotients of $F$, and their associated graded objects. If it
instead means all pro-prime quotients of the profinite group
$\widehat{\mathrm{GT}}_0$ itself, the proposed assertion is false:
the nonzero character $\Delta$ is already a quotient of that group
of order two, hence factors through its maximal pro-two quotient.
In that interpretation its value is determined by $V$, so the
image is a graph rather than $V\times\mathbf F_2$.

\paragraph{LB.7 (Refuted): the signs of LM8 B.31 are not the CM frame parities.}
The proposed identification in B.6(e) is false. The linked and
independently signed actual sources of LM8 B.31 coincide for
an elementary abelianization reason which concerns the powers
of torsion generators, whereas $\Delta$ concerns their
conjugating frames.

\emph{Proof and replacement.}
Let $C=\widehat{\mathrm{SL}_2(\mathbf Z)}$ with
$\rho^6=\sigma^4=1$, $\rho^3=\sigma^2$, and $e_0=\sigma\rho$.
Its abelianization is $\mathbf Z/12$, with $[\rho]=2$,
$[\sigma]=-3$, $[e_0]=-1$. An actual automorphism with
$u(e_0)=e_0^\mu$, $u(\rho)\sim\rho^{e_r}$ and
$u(\sigma)\sim\sigma^{e_s}$ therefore satisfies
\[
              e_r\equiv\mu\pmod6,\qquad
              e_s\equiv\mu\pmod4.
\]
The two exponent signs are forced by $\mu$; at multiplier one
they are both positive. The same calculation on the projective
modular group forces the corresponding positive exponent
classes there. Nevertheless the R2 witness has multiplier one
and CM frame parities $(1,0)$. Thus the two kinds of signs
cannot be identified. The strongest true statement is that
the exponent signs describe the cyclotomic power action,
while the CM parities are residual conjugator information
after that power action and the cusp shear have been fixed.
This addresses the proposed comparison with section 49.IV--V,
cote 140-3, scans 530, 549 and 559--560 as cited in LM8 B.31;
the original scans are not supplied here, so it makes no new
palaeographic assertion about them. \hfill$\square$
